% !TEX root = ../main_proposal.tex
% =====================================================================
%  BAB 2 — LANDASAN KEPUSTAKAAN
%  (istilah resmi Panduan Skripsi Filkom UB v3.0; bukan "Tinjauan Pustaka")
% =====================================================================

\chapter{\MakeUppercase{Landasan Kepustakaan}}

% ── 2.1 Kajian Pustaka ───────────────────────────────────────────────
\section{Kajian Pustaka}
\label{sec:2.1}

Kajian pustaka berisi kajian-kajian penelitian terdahulu yang bertujuan
untuk mendukung proses penelitian. Kajian pustaka dilakukan untuk
mendapatkan informasi yang relevan dengan topik penelitian ini. Tabel
\ref{tab:2.penelitian_terdahulu} memaparkan penelitian-penelitian yang
terkait dengan subjek yang diteliti.

{\small
\begin{longtable}{|p{0.17\textwidth}|p{0.18\textwidth}|p{0.30\textwidth}|p{0.19\textwidth}|}
  \caption{Daftar penelitian terdahulu}
  \label{tab:2.penelitian_terdahulu} \\
  \hline
  \centering\arraybackslash\textbf{Judul\newline Penelitian} & \centering\arraybackslash\textbf{Masalah} & \centering\arraybackslash\textbf{Metode} & \centering\arraybackslash\textbf{Hasil} \\
  \hline
  \endfirsthead
  \hline
  \centering\arraybackslash\textbf{Judul\newline Penelitian} & \centering\arraybackslash\textbf{Masalah} & \centering\arraybackslash\textbf{Metode} & \centering\arraybackslash\textbf{Hasil} \\
  \hline
  \endhead
  \hline
  \endfoot
  \hline
  \endlastfoot
    Development of a Speech Emotion Recognition Dataset for Indonesian
    (Santoso, Budianto and Dutono, 2026)
    &
    Dataset publik SER berbahasa Indonesia untuk penelitian masih sangat terbatas, dataset yang tersedia tidak mengakomodasi intonasi lokal, berskala kecil, terfragmentasi dan di rekam secara tidak natural.
    &
    Pembuatan dataset E-SERAVD melalui ekstraksi dan pembersihan audio dialog dari film komersial Indonesia, yang kemudian divalidasi ke dalam empat kelas emosi dasar oleh tiga penilai manusia. Dataset tersebut selanjutnya diproses menggunakan ekstraksi fitur MFCC dan dievaluasi kinerjanya menggunakan model CNN, tanpa melibatkan kontrol eksplisit terhadap sumber film
    pada pembagian data latih dan uji.
    &
    Dataset memiliki kualitas annotasi yang kuat. Model CNN dasar mencapai akurasi dan F1-score sebesar 86\% pada subset
    seimbang (1.000 sampel, 250/kelas), dengan emosi yg paling benar diprediksi adalah "Sedih"
    \\
    \hline
    End-to-End Transfer Learning for Speaker-Independent Cross-Language
    and Cross-Corpus Speech Emotion Recognition (Tang et al., 2023)
    &
    Evaluasi model SER yang tidak mengontrol kebocoran identitas
    penutur antara data latih dan uji berisiko melebih-lebihkan
    performa generalisasi model.
    &
    Skema evaluasi \emph{speaker-independent} yang memastikan penutur
    pada data uji tidak pernah muncul pada data latih. Skema ini
    diterapkan pada \emph{transfer learning end-to-end} lintas bahasa
    dan lintas korpus, dibandingkan dengan skema evaluasi konvensional
    yang tidak mengontrol identitas penutur.
    &
    Skema evaluasi \emph{speaker-independent} menghasilkan estimasi
    performa yang lebih andal dibandingkan skema yang tidak mengontrol
    identitas penutur.
    \\
    \hline
    Designing and Evaluating Speech Emotion Recognition Systems: A
    Reality Check Case Study with IEMOCAP (Antoniou et al., 2023)
    &
    Belum ada konsensus protokol evaluasi SER pada IEMOCAP, khususnya
    terkait kontrol identitas penutur antara data latih dan uji,
    sehingga performa antarpenelitian sulit dibandingkan secara adil.
    &
    Tinjauan kritis dan kompilasi hasil dari sejumlah publikasi SER pada
    IEMOCAP, membandingkan performa yang dilaporkan di bawah skema
    \emph{Speaker-Dependent} (SD, pembagian data acak) dengan skema
    \emph{Speaker-Independent} (SI, \emph{leave-one-speaker/session-out}).
    &
    Skema SD secara konsisten dilaporkan menghasilkan performa lebih
    tinggi dibandingkan skema SI pada metode yang sebanding; satu model
    HuBERT pada penelitian yang sama melaporkan 79,58\% (SD) berbanding
    73,01\% (SI) pada metrik \emph{weighted accuracy}.
    \\
    \hline
    Evaluating Self-Supervised Speech Representations for Speech
    Emotion Recognition (Atmaja and Sasou, 2022)
    &
    Belum diketahui apakah keunggulan representasi \emph{self-supervised}
    atas fitur \emph{hand-crafted} pada SER berlaku universal lintas model dan
    bahasa.
    &
    Evaluasi sistematis 19 representasi \emph{self-supervised} dan 1
    fitur \emph{hand-crafted} (filterbank) pada 5 dataset SER berbahasa Inggris dan
    Jepang. Setiap representasi diuji dengan \emph{classifier
    feed-forward} yang identik, kemudian signifikansi antarrepresentasi
    diuji memakai uji Friedman dan dihitung \emph{effect size}-nya.
    &
    Representasi \emph{self-supervised} unggul signifikan (filterbank
    rank terakhir dari 20), namun XLSR-53 spesifik kalah dari model
    monolingual Inggris; tidak ada dataset Indonesia diuji.
    \\
    \hline
    Advancing End-to-End Indonesian Speech Emotion Recognition Using
    XLSR-Wav2Vec2 (Nugroho et al., 2025)
    &
    Performa SER berbahasa Indonesia belum diketahui secara komparatif
    antara metode klasik dan representasi \emph{self-supervised}
    multibahasa.
    &
    Perbandingan SVM, CNN, dan dua model dasar XLSR-Wav2Vec2 yang telah
    di-\emph{fine-tune} untuk ucapan Indonesia menggunakan data Mozilla
    Common Voice serta OpenSLR berbahasa Jawa dan Sunda. Keempat model
    diuji pada dataset IndoWaveSentiment untuk membandingkan performa
    metode klasik dan representasi \emph{self-supervised} multibahasa.
    &
    Model dasar XLSR-Wav2Vec2 yang telah disesuaikan untuk ucapan
    Indonesia mencapai akurasi tertinggi sebesar 91,67\%, mengungguli
    metode klasik.
    \\
\end{longtable}
}

% ── 2.2 Speech Emotion Recognition ───────────────────────────────────
\section{Speech Emotion Recognition (SER)}
\label{sec:2.2}

\emph{Speech emotion recognition} (SER) adalah bidang kajian yang
berupaya mengenali kondisi emosi seseorang secara otomatis dari sinyal
ucapannya, dengan memanfaatkan ciri-ciri akustik seperti intonasi, nada
dasar (\emph{pitch}), energi, dan ritme bicara yang secara empiris
berkorelasi dengan keadaan emosional penutur (Dar and Delhibabu, 2024).
Gambar \ref{fig:2.ser} merangkum cakupan bidang SER dari tiga sisi
utama, yaitu sumber data, jenis fitur, dan jenis \emph{classifier} yang
digunakan. Dari sisi sumber data, rekaman ucapan emosional dapat
diperoleh secara \emph{acted} (diperankan aktor mengikuti skenario
emosi tertentu), \emph{elicited} (dipancing muncul secara spontan
melalui stimulus atau situasi tertentu), maupun \emph{natural} (diambil
langsung dari situasi nyata tanpa rekayasa, misalnya dari dialog film
seperti yang digunakan pada penelitian ini). Dari sisi fitur,
representasi yang diekstraksi dari sinyal ucapan terbagi menjadi fitur
akustik (\emph{acoustic}, mis. \emph{prosodic}, spektral, dan kualitas
suara) dan fitur non-akustik (\emph{non-acoustic}, mis. linguistik,
wacana, serta ekspresi wajah dan gestur pada modalitas multimodal). Dari
sisi \emph{classifier}, metode yang digunakan terbagi menjadi metode
pembelajaran mesin konvensional (\emph{machine learning based}, mis.
SVM, HMM, KNN) dan metode pembelajaran mendalam (\emph{deep learning
based}, mis. CNN, RNN, LSTM). Penelitian ini berada pada irisan data
\emph{natural} berbasis film, fitur akustik, dan \emph{classifier}
berbasis pembelajaran mendalam.

\begin{figure}[H]
  \centering
  \gambarbesar{\includegraphics{SER-Overview-Der-Delbabu-2024.png}}
  \caption{Cakupan basis data, fitur, dan \emph{classifier} pada SER}
  \label{fig:2.ser}
  \normalfont Sumber : Dar and Delhibabu (2024)
\end{figure}

Cakupan penerapan SER terus meluas, mulai dari asisten virtual yang
responsif terhadap keadaan pengguna, layanan pelanggan otomatis yang
dapat mendeteksi frustrasi penelepon, hingga sistem pemantauan kesehatan
mental yang menandai perubahan pola emosi dari waktu ke waktu.
Perkembangan representasi \emph{self-supervised} seperti Wav2Vec2 dan
XLSR-53 menjadi salah satu pendorong utama peningkatan akurasi SER
dibandingkan pendekatan fitur \emph{hand-crafted} konvensional.

% ── 2.3 Data Leakage / Group Leakage ──────────────────────────────────
\section{Data Leakage / Group Leakage}
\label{sec:2.3}

Kebocoran data (\emph{data leakage}) (Kapoor and Narayanan, 2023) adalah
kondisi ketika informasi dari data uji secara tidak sengaja ikut
memengaruhi proses pelatihan model, sehingga performa yang terukur pada
data uji tidak lagi mencerminkan kemampuan generalisasi model yang
sesungguhnya pada data baru. Salah satu bentuk kebocoran data yang
sering luput diperhatikan adalah \emph{group leakage}, yaitu kebocoran
yang terjadi ketika beberapa sampel data memiliki keterkaitan atau
berasal dari satu sumber/kelompok yang sama (dikategorikan sebagai
\emph{nonindependence} antara sampel latih dan uji oleh Kapoor and
Narayanan, 2023) -- misalnya beberapa klip ucapan yang berasal dari
rekaman film yang sama -- namun sampel-sampel berkelompok tersebut
tersebar ke sisi latih dan uji sekaligus akibat pembagian data yang
dilakukan secara acak murni terhadap satuan sampel individual (skema
\emph{random}). Ketika hal ini terjadi, model berpotensi mempelajari
ciri-ciri spesifik kelompok tersebut -- misalnya gaya produksi,
karakteristik rekaman, atau suara latar khas suatu film -- sebagai
jalan pintas untuk memprediksi label, alih-alih benar-benar mempelajari
pola yang ingin dikenali (dalam konteks penelitian ini, pola akustik
emosi). Akibatnya, performa yang dilaporkan pada skema \emph{random}
cenderung terlihat lebih tinggi daripada performa sesungguhnya yang
akan diperoleh model pada data dari sumber yang benar-benar baru.

Skema pembagian data \emph{movie-independent} (satu ragam
\emph{grouped split}) mengatasi risiko \emph{group leakage} tersebut
dengan menjaga agar seluruh klip yang berasal dari film yang sama
selalu berada pada sisi \emph{fold} yang sama -- tidak pernah ada satu
film yang klipnya tersebar ke sisi latih dan uji sekaligus. Dengan
demikian, ketika suatu film digunakan sebagai data uji, model
benar-benar belum pernah ``melihat'' karakteristik rekaman film
tersebut selama pelatihan, sehingga performa yang terukur mencerminkan
kemampuan generalisasi model pada sumber rekaman yang genuinely baru.
Preseden penerapan skema \emph{grouped} semacam ini pada SER dilaporkan
oleh Tang et al. (2023), yang menerapkan skema evaluasi
\emph{speaker-independent} -- varian \emph{grouped split} berdasarkan
identitas penutur, bukan berdasarkan film -- untuk mencegah model
mempelajari isyarat spesifik penutur alih-alih pola emosi yang
genuinely general. Penelitian ini menerapkan prinsip yang sama pada
satuan film, karena unit kebocoran yang relevan pada dataset yang
digunakan adalah sumber film, bukan identitas penutur individual.

Penelitian ini membandingkan performa algoritma yang diuji di bawah
kedua skema tersebut -- \emph{random} dan \emph{movie-independent} --
untuk mengukur secara langsung seberapa besar performa yang dilaporkan
berpotensi terinflasi oleh \emph{group leakage} berbasis film, dan
apakah besar inflasi tersebut berbeda antara algoritma berbasis fitur
\emph{hand-crafted} dan algoritma berbasis representasi \emph{self-supervised}.

% ── 2.4 Fitur Hand-Crafted ────────────────────────────────────────────
\section{Fitur Hand-Crafted}
\label{sec:2.4}

Fitur \emph{hand-crafted} adalah representasi sinyal yang diekstraksi
secara manual melalui rangkaian transformasi sinyal yang telah
ditentukan terlebih dahulu (\emph{fixed}) berdasarkan teori dan
pengetahuan domain, alih-alih dipelajari otomatis oleh model dari data.
Pada domain pemrosesan ucapan, fitur \emph{hand-crafted} umumnya
dirancang untuk menangkap karakteristik akustik yang secara teoretis
relevan dengan tugas yang ingin diselesaikan, misalnya intonasi,
energi, nada dasar (\emph{pitch}), atau distribusi frekuensi sinyal.
Karena rangkaian transformasinya sudah ditentukan sejak awal, fitur
\emph{hand-crafted} tidak memerlukan data berlabel maupun proses
pembelajaran untuk diperoleh, berbeda dengan representasi yang
dipelajari model dari data (\emph{learned representations}) yang
parameternya disesuaikan melalui proses pelatihan.

Pendekatan fitur \emph{hand-crafted} banyak digunakan pada penelitian
SER sebelum berkembangnya representasi yang dipelajari model, karena
tidak memerlukan data pra-latih berskala besar dan relatif murah secara
komputasi untuk diekstraksi. Salah satu fitur \emph{hand-crafted} yang
paling umum digunakan pada pemrosesan ucapan adalah Mel-Frequency
Cepstral Coefficients (MFCC), yang digunakan sebagai representasi fitur
\emph{hand-crafted} pada penelitian ini.

% ── 2.5 Mel-Frequency Cepstral Coefficients (MFCC) ───────────────────
\section{Mel-Frequency Cepstral Coefficients (MFCC)}
\label{sec:2.5}

Mel-Frequency Cepstral Coefficients (MFCC) adalah representasi fitur
akustik \emph{hand-crafted} yang secara luas digunakan pada tugas pemrosesan ucapan,
dirancang untuk meniru cara sistem pendengaran manusia mempersepsikan
frekuensi bunyi secara tidak linear (Davis and Mermelstein, 1980).
Alih-alih menggunakan skala frekuensi linear (Hz), MFCC memetakan
spektrum sinyal ke skala mel, yang lebih sensitif terhadap perbedaan
pada frekuensi rendah dibandingkan frekuensi tinggi, sebagaimana
Persamaan \ref{eq:2.mel}.

\begin{equation}
  m = 2595 \log_{10}\left(1 + \frac{f}{700}\right)
  \label{eq:2.mel}
\end{equation}

\noindent
di mana $f$ adalah frekuensi dalam Hz dan $m$ adalah frekuensi
terskala mel yang bersesuaian.

Ekstraksi MFCC dari sinyal ucapan mentah dilakukan melalui serangkaian
tahap sebagaimana digambarkan pada Gambar \ref{fig:2.mfcc}. Sinyal
terlebih dahulu melalui \emph{pre-emphasis} untuk menguatkan komponen
frekuensi tinggi, kemudian dibagi menjadi bingkai-bingkai pendek
(\emph{framing}) yang tumpang tindih dan dikalikan fungsi jendela
(\emph{windowing}) untuk mengurangi efek diskontinuitas di tepi bingkai.
Setiap bingkai kemudian ditransformasi ke domain frekuensi menggunakan
\emph{Fast Fourier Transform} (FFT) untuk memperoleh spektrum daya, yang
selanjutnya difilter menggunakan sekumpulan filter \emph{bank} berskala
mel (Persamaan \ref{eq:2.mel}) untuk memperoleh energi pada tiap pita
frekuensi mel. Energi tersebut diambil logaritmanya untuk mendekati
persepsi kenyaringan manusia yang bersifat logaritmik, kemudian
ditransformasi menggunakan \emph{Discrete Cosine Transform} (DCT) untuk
mendekorelasi energi antarpita dan menghasilkan koefisien cepstral
akhir.

\begin{figure}[H]
  \centering
  \gambarbesar{\includegraphics{MFCC-Block-DIagram.png}}
  \caption{Pipeline ekstraksi Mel-Frequency Cepstral Coefficients (MFCC)}
  \label{fig:2.mfcc}
  \normalfont Sumber : Alashban and Alotaibi (2023)
\end{figure}

Berbeda dengan representasi yang dipelajari model dari data
(\emph{learned representations}), MFCC dihitung langsung dari sinyal
melalui rangkaian transformasi sinyal yang telah ditentukan
(\emph{fixed}), tanpa proses pembelajaran. Karakteristik inilah yang
mendasari penggunaan MFCC sebagai fitur \emph{hand-crafted} pada
algoritma \emph{baseline} dalam penelitian ini.

% ── 2.6 Convolutional Neural Network (CNN) ───────────────────────────
\section{Convolutional Neural Network (CNN)}
\label{sec:2.6}

\emph{Convolutional neural network} (CNN) adalah arsitektur jaringan
saraf tiruan yang dirancang untuk memproses data yang memiliki struktur
grid, seperti citra dua dimensi atau spektrogram ucapan, dengan
memanfaatkan operasi konvolusi untuk mengekstraksi pola-pola lokal
secara bertingkat (LeCun et al., 1998). Arsitektur CNN umumnya terdiri
atas tiga jenis lapisan utama, sebagaimana digambarkan pada Gambar
\ref{fig:2.cnn}: (1) lapisan konvolusi (\emph{convolutional layer}),
yang menerapkan sejumlah filter/\emph{kernel} berukuran kecil secara
bergeser di sepanjang data masukan untuk mendeteksi pola lokal seperti
tepi atau tekstur pada tahap awal, dan pola yang lebih kompleks pada
lapisan berikutnya; (2) lapisan \emph{pooling}, yang mereduksi dimensi
spasial representasi sambil mempertahankan informasi penting, sehingga
jaringan lebih tahan terhadap pergeseran kecil pada data masukan dan
lebih efisien secara komputasi; dan (3) lapisan terhubung penuh
(\emph{fully-connected layer}), yang menggabungkan seluruh fitur yang
telah diekstraksi untuk menghasilkan keluaran akhir, misalnya
probabilitas kelas pada tugas klasifikasi. Secara garis besar, lapisan
konvolusi dan \emph{pooling} membentuk tahap ekstraksi fitur
(\emph{feature extraction}), sedangkan lapisan terhubung penuh membentuk
tahap klasifikasi (\emph{classification}), sebagaimana ditunjukkan pada
Gambar \ref{fig:2.cnn}.

\begin{figure}[H]
  \centering
  \gambarbesar{\includegraphics{Basic-CNN-Diagram.png}}
  \caption{Arsitektur umum Convolutional Neural Network (CNN)}
  \label{fig:2.cnn}
  \normalfont Sumber : Phung and Rhee (2019)
\end{figure}

Kemampuan lapisan konvolusi mendeteksi pola lokal menjadikan CNN cocok
untuk memproses fitur MFCC, yang berbentuk grid dua dimensi berupa
koefisien cepstral pada sumbu vertikal dan bingkai waktu pada sumbu
horizontal -- pola akustik emosi yang relevan seringkali muncul sebagai
struktur lokal pada grid tersebut, misalnya perubahan energi pada
rentang waktu dan pita frekuensi tertentu. Karakteristik inilah yang
mendasari penggunaan CNN sebagai \emph{classifier} pada algoritma
\emph{baseline} dalam penelitian ini, mengikuti pendekatan yang sudah
banyak digunakan pada penelitian SER berbasis fitur \emph{hand-crafted}.

% ── 2.7 Pembelajaran Self-Supervised ─────────────────────────────────
\section{Pembelajaran Self-Supervised}
\label{sec:2.7}

Pembelajaran \emph{self-supervised} adalah paradigma pembelajaran mesin
yang melatih model pada data tanpa label dengan cara menciptakan sinyal
pengawasan (\emph{supervision signal}) dari data itu sendiri, alih-alih
memerlukan label yang dianotasi manusia (Liu et al., 2023). Model
dilatih untuk menyelesaikan \emph{pretext task}, yaitu tugas bantu yang
dirancang agar penyelesaiannya mengharuskan model mempelajari struktur
dan pola yang mendasari data, misalnya memprediksi bagian data yang
disembunyikan dari bagian yang tersisa. Sebagaimana ditunjukkan pada
Gambar \ref{fig:2.ssl}, bobot \emph{encoder} yang telah dilatih pada
\emph{pretext task} menggunakan data tanpa label kemudian ditransfer
(\emph{transfer weights}) sebagai titik awal model pada tugas hilir
(\emph{downstream task}) yang sesungguhnya, sebelum diadaptasi lebih
lanjut (\emph{fine-tuned}) menggunakan data berlabel, seperti
klasifikasi emosi pada penelitian ini.

\begin{figure}[H]
  \centering
  \gambarbesar{\includegraphics{Self-Supervised-Diagram.png}}
  \caption{Alur umum pembelajaran Self-Supervised}
  \label{fig:2.ssl}
  \normalfont Sumber : Del Pup and Atzori (2023)
\end{figure}

Keunggulan utama paradigma ini adalah kemampuannya memanfaatkan data
tanpa label dalam jumlah besar, yang jauh lebih mudah dan murah diperoleh
dibandingkan data berlabel. Karakteristik inilah yang mendasari
penggunaannya secara luas pada domain ucapan, termasuk pada backbone
Wav2Vec2 dan XLSR-53 yang digunakan dalam penelitian ini.

% ── 2.8 Wav2Vec2 ──────────────────────────────────────────────────────
\section{Wav2Vec2}
\label{sec:2.8}

Wav2Vec2 adalah arsitektur pembelajaran \emph{self-supervised} untuk
representasi ucapan yang dilatih pada volume data ucapan tanpa label
berskala besar (Baevski et al., 2020). Arsitekturnya terdiri atas tiga
komponen utama, sebagaimana digambarkan pada Gambar \ref{fig:2.wav2vec2}:
(1) \emph{feature encoder} berbasis \emph{convolutional neural network}
(CNN) yang mengubah gelombang ucapan mentah menjadi representasi laten
$Z$ pada resolusi waktu yang lebih rendah; (2) \emph{context network}
berbasis Transformer yang mengolah $Z$ (dengan sebagian di antaranya
disembunyikan/\emph{masked}) menjadi representasi terkontekstualisasi
$C$; dan (3) modul kuantisasi yang mendiskretkan $Z$ menjadi target
kuantisasi $Q$, khusus digunakan pada tahap pra-latih untuk menghitung
\emph{contrastive loss} yang mendorong model membedakan representasi
kuantisasi yang benar dari sejumlah pengganggu (\emph{distractors}).

\begin{figure}[H]
  \centering
  \gambarbesar{\includegraphics{WAV2VEC2-Diagram.png}}
  \caption{Arsitektur Wav2Vec2}
  \label{fig:2.wav2vec2}
  \normalfont Sumber : Ovalle et al. (2024)
\end{figure}

Representasi $C$ yang dihasilkan oleh \emph{context network} inilah yang
menjadi keluaran utama Wav2Vec2 untuk dimanfaatkan pada tugas hilir
(\emph{downstream task}), termasuk SER pada penelitian ini.

% ── 2.9 XLSR-53 ───────────────────────────────────────────────────────
\section{XLSR-53}
\label{sec:2.9}

XLSR merupakan varian multibahasa dari Wav2Vec2 yang dilatih pada data
ucapan tanpa label dari banyak bahasa sekaligus, menggunakan arsitektur
dan tujuan pra-latih yang sama, namun dengan korpus pra-latih gabungan
lintas bahasa (Conneau et al., 2021). Karena dilatih lintas bahasa,
representasi yang dihasilkan XLSR cenderung menangkap struktur akustik
yang lebih umum dan dapat menggeneralisasi lintas bahasa dibandingkan
model yang hanya dilatih pada satu bahasa saja. XLSR-53 adalah varian
XLSR yang dilatih pada korpus gabungan dari 53 bahasa. Karakteristik
inilah yang mendasari pemilihan XLSR-53 sebagai backbone algoritma SOTA
dalam penelitian ini: bobot dasar backbone tetap dibekukan
(\emph{frozen}), namun diberi matriks adaptasi tambahan berperingkat
rendah (LoRA) yang dilatih untuk menyesuaikan representasi ke tugas
SER, alih-alih memperbarui seluruh parameter backbone secara langsung.

% ── 2.10 Low-Rank Adaptation (LoRA) ──────────────────────────────────
\section{Low-Rank Adaptation (LoRA)}
\label{sec:2.10}

\emph{Low-Rank Adaptation} (LoRA) adalah strategi \emph{fine-tuning}
efisien-parameter yang membekukan seluruh bobot model dasar $W_0$ dan
hanya melatih sepasang matriks tambahan berperingkat rendah
(\emph{low-rank}) yang disisipkan secara paralel (Hu et al., 2022).
Untuk suatu lapisan dengan bobot pra-latih $W_0 \in \mathbb{R}^{d \times
k}$, LoRA merepresentasikan pembaruan bobot $\Delta W$ sebagai perkalian
dua matriks berperingkat rendah $B \in \mathbb{R}^{d \times r}$ dan $A
\in \mathbb{R}^{r \times k}$, dengan $r \ll \min(d, k)$, sehingga keluaran
lapisan tersebut untuk masukan $x$ dihitung sebagaimana Persamaan
\ref{eq:2.lora}.

\begin{equation}
  h = W_0 x + \Delta W x = W_0 x + BA x
  \label{eq:2.lora}
\end{equation}

\noindent
Gambar \ref{fig:2.lora} membandingkan skema tersebut dengan \emph{full
fine-tuning} konvensional. Pada \emph{full fine-tuning} (kiri),
pembaruan bobot $\Delta W$ dilatih penuh sebagai satu matriks padat
berukuran sama dengan $W_0$. Pada LoRA (kanan), bobot dasar $W_0$
dibekukan (ditandai ikon salju) sementara hanya matriks $A$ dan $B$
berperingkat rendah yang dilatih (ditandai ikon api), sebelum hasil
keduanya dijumlahkan menjadi keluaran $h$.

\begin{figure}[H]
  \centering
  \gambarbesar{\includegraphics{Regular-Fine-TUne-VS-LoRA-diagram-.png}}
  \caption{Perbandingan full fine-tuning dan LoRA}
  \label{fig:2.lora}
  \normalfont Sumber : Lu et al. (2024)
\end{figure}

Karena jumlah parameter pada $A$ dan $B$ jauh lebih sedikit daripada
jumlah parameter $W_0$, LoRA menawarkan efisiensi komputasi dan memori
yang signifikan dibandingkan \emph{full fine-tuning}, yang memperbarui
seluruh parameter $W_0$ secara langsung. Efisiensi ini relevan bagi
penelitian dengan keterbatasan sumber daya komputasi. Efektivitas LoRA
pada arsitektur berbasis Wav2Vec2 secara spesifik telah dibuktikan pada
tugas deteksi audio sintetis, di mana LoRA mencapai performa setara
dengan \emph{full fine-tuning} sekaligus mengurangi jumlah parameter
yang dilatih hingga 198 kali lipat (Wang et al., 2023). Pada tugas SER
secara lebih spesifik, kombinasi LoRA dengan metode
\emph{parameter-efficient fine-tuning} lainnya terbukti melampaui
performa \emph{full fine-tuning} pada backbone Wav2Vec2 dan HuBERT,
sekaligus hanya melatih sekitar 2--3\% dari total parameter
(Lashkarashvili et al., 2024). Karakteristik inilah yang mendasari
penggunaan LoRA untuk mengadaptasi backbone XLSR-53 pada algoritma SOTA
dalam penelitian ini, alih-alih memperbarui seluruh parameter backbone
secara langsung.

% ── 2.11 Evaluasi Model ──────────────────────────────────────────────
\section{Evaluasi Model}
\label{sec:2.11}

Akurasi, presisi, recall, dan F1-score adalah empat metrik yang umum
digunakan untuk mengevaluasi performa model klasifikasi (Powers, 2011).
Akurasi mengukur proporsi prediksi yang benar terhadap keseluruhan data
uji, sehingga memberikan gambaran umum performa model, sebagaimana
Persamaan \ref{eq:2.akurasi}.

\begin{equation}
  \text{Akurasi} = \frac{TP + TN}{TP + TN + FP + FN}
  \label{eq:2.akurasi}
\end{equation}

\noindent
di mana $TP$ = \emph{True Positive}, $TN$ = \emph{True Negative}, $FP$ =
\emph{False Positive}, $FN$ = \emph{False Negative}.

Akan tetapi, pada data dengan sebaran kelas yang tidak seimbang, akurasi
dapat memberikan gambaran yang menyesatkan, karena model yang selalu
memprediksi kelas mayoritas tetap dapat memperoleh akurasi tinggi tanpa
benar-benar mengenali kelas minoritas. Presisi mengukur proporsi
prediksi positif yang benar-benar tepat, sebagaimana Persamaan
\ref{eq:2.presisi}.

\begin{equation}
  \text{Presisi} = \frac{TP}{TP + FP}
  \label{eq:2.presisi}
\end{equation}

\noindent
Recall mengukur proporsi kasus positif yang berhasil dikenali model,
sebagaimana Persamaan \ref{eq:2.recall}.

\begin{equation}
  \text{Recall} = \frac{TP}{TP + FN}
  \label{eq:2.recall}
\end{equation}

\noindent
Presisi dan recall sering kali saling bertukar (\emph{trade-off}):
menaikkan salah satu cenderung menurunkan yang lain. F1-score mengatasi
keterbatasan ini dengan merepresentasikan rata-rata harmonis keduanya,
sebagaimana Persamaan \ref{eq:2.f1}.

\begin{equation}
  F_1 = 2 \cdot \frac{\text{Presisi} \times \text{Recall}}{\text{Presisi} + \text{Recall}}
  \label{eq:2.f1}
\end{equation}

Keempat metrik ini digunakan bersama dalam penelitian ini untuk
memberikan gambaran performa yang lebih menyeluruh terhadap setiap
kondisi yang dibandingkan. Meskipun skema kelas emosi pada dataset yang
digunakan dalam penelitian ini seimbang antarkelas, akurasi saja tidak
menunjukkan pola kesalahan per kelas, sedangkan presisi, recall, dan
F1-score memberikan gambaran lebih rinci mengenai kelas emosi mana yang
lebih sering tertukar atau lebih sulit dikenali model.

% ── 2.12 Validasi Silang K-Fold ───────────────────────────────────────
\section{Validasi Silang K-Fold}
\label{sec:2.12}

Validasi silang k-fold adalah teknik evaluasi yang membagi data menjadi
$k$ bagian (\emph{fold}) berukuran relatif sama, dengan setiap bagian
secara bergantian digunakan sebagai data uji sementara $k-1$ bagian
lainnya digunakan sebagai data latih, sehingga setiap bagian data pernah
menjadi data uji tepat satu kali (Kohavi, 1995). Gambar
\ref{fig:2.kfold} mengilustrasikan mekanisme umum tersebut: \emph{fold}
uji (biru) berpindah bergiliran pada tiap iterasi sementara \emph{fold}
sisanya berperan sebagai data latih, dengan hasil evaluasi $E_i$ dari
tiap iterasi dirata-ratakan menjadi estimasi performa akhir $E$.
Penelitian ini secara spesifik menggunakan $k=6$.

\begin{figure}[H]
  \centering
  \gambarbesar{\includegraphics{K-Fold-Diagram.png}}
  \caption{Ilustrasi mekanisme validasi silang k-fold}
  \label{fig:2.kfold}
  \normalfont Sumber : Ashfaque and Iqbal (2019)
\end{figure}

Teknik ini menghasilkan $k$ observasi akurasi dan F1-score, satu untuk
setiap fold, sehingga memungkinkan pengujian konsistensi performa
antar-fold untuk setiap kondisi yang diuji, bukan hanya satu nilai
tunggal dari satu kali pengujian.

% ── 2.13 Pengujian Signifikansi Statistik ────────────────────────────
\section{Pengujian Signifikansi Statistik}
\label{sec:2.13}

Untuk menguji apakah perbedaan performa antarkondisi hasil validasi
silang k-fold bersifat signifikan secara statistik, $k$ observasi yang
diperoleh terlebih dahulu diuji normalitasnya menggunakan uji
Shapiro-Wilk, yang menghitung statistik $W$ sebagaimana Persamaan
\ref{eq:2.shapirowilk} (Shapiro and Wilk, 1965).

\begin{equation}
  W = \frac{\left(\sum_{i=1}^{n} a_i x_{(i)}\right)^2}{\sum_{i=1}^{n} (x_i - \bar{x})^2}
  \label{eq:2.shapirowilk}
\end{equation}

\noindent
di mana $x_{(i)}$ adalah nilai observasi ke-$i$ setelah diurutkan
menaik (statistik urutan), $\bar{x}$ adalah rata-rata observasi, $n$
adalah jumlah observasi (jumlah fold), dan $a_i$ adalah konstanta yang
diturunkan dari rata-rata, ragam, dan kovarians statistik urutan sampel
berdistribusi normal berukuran $n$. Nilai $W$ yang mendekati 1
mengindikasikan data mendekati distribusi normal.

Apabila data berdistribusi normal, pengujian dilanjutkan dengan
\emph{paired t-test}, yang menghitung statistik uji $t$ dari selisih
berpasangan $d_i$ antara dua kondisi pada fold yang sama, sebagaimana
Persamaan \ref{eq:2.ttest}.

\begin{equation}
  t = \frac{\bar{d}}{s_d / \sqrt{n}}
  \label{eq:2.ttest}
\end{equation}

\noindent
di mana $\bar{d}$ adalah rata-rata selisih berpasangan, $s_d$ adalah
simpangan baku selisih tersebut, dan $n$ adalah jumlah fold. Apabila
data tidak berdistribusi normal, digunakan \emph{Wilcoxon signed-rank
test} sebagai alternatif non-parametrik yang tidak mengasumsikan
distribusi tertentu, dengan statistik uji $W$ sebagaimana Persamaan
\ref{eq:2.wilcoxon} (Wilcoxon, 1945).

\begin{equation}
  W = \sum_{i=1}^{n} \operatorname{sgn}(d_i) \, R_i
  \label{eq:2.wilcoxon}
\end{equation}

\noindent
di mana $d_i$ adalah selisih berpasangan pada fold ke-$i$, $R_i$ adalah
peringkat (\emph{rank}) dari $|d_i|$ di antara seluruh selisih mutlak
$|d_1|, \dots, |d_n|$, dan $\operatorname{sgn}(d_i)$ bernilai $+1$
apabila $d_i>0$ dan $-1$ apabila $d_i<0$.
